% ============================================================
% Chapter 6: Experimental Design
%   Section: Construction of the Custom Validation Set
%     Subsection: Verification of the Labels
% Included from chapters/06-experimental-design/03-custom-validation-set/custom-validation-set.tex
% ============================================================

\subsection{Verification of the Labels}
\label{subsec:custom-review}

The set is labelled by a single annotator, so inter-annotator
agreement cannot be measured. What can be measured is
\textbf{intra-annotator agreement}: whether the guide, applied by the
same person at a different time, produces the same label. Re-evaluating
a sample is a requirement of the protocol, not an option: without it
there is no way to tell a consistent label from one that would have come
out differently on another day. Once labelling is complete, 20\% of the
\emph{test half} (Subsection~\ref{subsec:custom-split}) is labelled a
second time --- the claims the reported numbers rest on:

\begin{itemize}
    \item \textbf{The sample is not chosen by the annotator.} It is
    the first $\lceil 0.2\,n \rceil$ records of the test half in the
    order of a cryptographic hash (SHA-256) of their identifier, which
    is fixed in advance and independent of content: 14 of the 70
    test claims, listed in the split's record
    (\texttt{split.json}, \texttt{recheck}).
    \item \textbf{The second pass is blind.} The review view shows
    the claim, its source article, dates, topic and evidence links,
    but not the first label, the \texttt{onlyOwnSource} flag or the
    annotator's note, all of which reveal it.
    \item \textbf{It takes place at least a week after the first
    pass}, so that the second label is produced by the guide rather
    than recalled.
\end{itemize}

Agreement is reported as observed agreement and as Cohen's $\kappa$
\cite{cohen1960coefficient}, which discounts the agreement two
labellings would reach by chance given how often each used every
label; this matters here because the label distribution is expected
to be uneven. Following the conventions reviewed by Artstein and
Poesio \cite{artstein2008intercoder}, $\kappa \geq 0.6$ (``substantial''
on the scale of Landis and Koch \cite{landis1977measurement}) is taken
as the threshold for the guide to be considered consistently
applicable; below it, the rules behind the disagreements are revised
and the affected strata relabelled.

Disagreements are resolved against the guide and, where the first
label proves wrong, the record is corrected. Agreement remains
measured on the \emph{first} label, which the record keeps beside the
second: correcting a record after its review must not turn a
disagreement into an agreement.

\paragraph{The same requirement, applied to the system.} A label can
be inconsistent; so can a verdict. The pipeline searches a web that
changes between two runs, through engines that answer differently from
one request to the next, and a local model at temperature zero still
varies. The same 14 claims are therefore also evaluated a second time
by the system, with the same configuration, and the two runs are
compared verdict against verdict (Subsection~\ref{subsec:consistency}).
The annotator's self-agreement and the system's run-to-run agreement
are both reported as Cohen's $\kappa$ on the same claims, so the two
consistencies can be read side by side.
